\documentclass[preprint,12pt]{elsarticle}

\usepackage{graphicx}
\usepackage{booktabs}
\usepackage{amsmath}
\usepackage{algorithm}
\usepackage{algorithmic}
\usepackage{hyperref}
\usepackage{multirow}
\usepackage{xcolor}
\usepackage{listings}
\usepackage{tcolorbox}
\usepackage{enumitem}

% Colors for prompt boxes
\definecolor{rolecolor}{RGB}{230,240,250}
\definecolor{goalcolor}{RGB}{240,250,230}
\definecolor{rulecolor}{RGB}{255,250,230}
\definecolor{outputcolor}{RGB}{250,240,245}

\journal{Engineering Applications of Artificial Intelligence}

\begin{document}

\begin{frontmatter}

\title{An Agentic AI Framework for Automated Process Improvement through LLM-Guided Hypothesis Generation and Dual-Simulator Validation}


\begin{abstract}
Process improvement remains a critical challenge for organizations seeking operational excellence, yet traditional approaches require extensive manual analysis and domain expertise. We present an agentic AI framework that automates the discovery and validation of process improvements by combining Large Language Model (LLM)-based hypothesis generation with dual-simulator validation. Our approach extracts baseline simulation models from event logs using established process mining techniques, then iteratively generates improvement hypotheses targeting actionable resource-related interventions. Each hypothesis is validated through parallel execution of Discrete Event Simulation (DES) and Agent-Based Simulation (ABS), with averaged results providing robust validation that neither simulator alone can achieve. The framework supports three intervention types: responsibility shifting between roles, resource allocation adjustments, and working calendar modifications. We evaluate our approach through two complementary studies: (1) a quantitative benchmark comparing LLM-guided search against systematic grid search across 15 fault-injection experiments, demonstrating significant speedup in convergence; and (2) a qualitative study with domain experts validating the actionability of recommendations on real-world process logs. Our results show that the agentic approach consistently identifies valid improvements with fewer iterations than exhaustive search while generating recommendations that practitioners find actionable and useful.
\end{abstract}

\begin{keyword}
Agentic AI \sep Process Improvement \sep Large Language Models \sep Business Process Simulation \sep Digital Twins \sep Process Mining
\end{keyword}

\end{frontmatter}

%% ============================================================================
\section{Introduction}
\label{sec:intro}
%% ============================================================================

Business process improvement is a fundamental concern for organizations striving to enhance operational efficiency, reduce costs, and improve customer satisfaction. Traditional process improvement methodologies rely heavily on manual analysis by domain experts who examine process models, identify bottlenecks, and propose interventions based on experience and intuition. This approach, while valuable, is time-consuming, subjective, and often fails to explore the full space of possible improvements.

The emergence of process mining has enabled data-driven analysis of business processes through event logs captured by information systems~\cite{van2016process}. Process mining techniques can automatically discover process models, check conformance, and analyze performance. However, the transition from descriptive analytics (``what happened'') to prescriptive analytics (``what should we do'') remains challenging. While process mining can identify bottlenecks and inefficiencies, recommending specific, actionable improvements requires understanding complex interdependencies and predicting the impact of proposed changes.

Recent advances in Large Language Models (LLMs) have demonstrated remarkable capabilities in reasoning, planning, and decision-making across diverse domains. The emergence of \textit{agentic AI}---systems that can perceive, reason, plan, and act autonomously---opens new possibilities for intelligent automation in industrial settings. In the context of process improvement, LLMs can leverage their broad knowledge to generate creative improvement hypotheses while simulation provides rigorous validation of predicted impacts.

In this paper, we present an agentic AI framework for automated process improvement that combines LLM-based hypothesis generation with dual-simulator validation. Our approach addresses a key limitation of both pure LLM-based recommendations (which may be unrealistic or unvalidated) and pure simulation-based optimization (which requires predefined search spaces and may miss creative solutions). The framework operates as follows:

\begin{enumerate}
    \item \textbf{Baseline Discovery}: Extract a simulation model from historical event logs using established process mining and simulation model discovery techniques.
    \item \textbf{Iterative Improvement}: Generate improvement hypotheses using an LLM, guided by baseline performance metrics and previous attempts.
    \item \textbf{Dual Validation}: Validate each hypothesis through parallel execution of Discrete Event Simulation (DES) and Agent-Based Simulation (ABS), averaging results for robust assessment.
    \item \textbf{Report Generation}: Produce actionable improvement recommendations with quantified expected impacts.
\end{enumerate}

Our framework focuses on \textit{actionable} interventions---changes that organizations can realistically implement. Rather than simply identifying ``reduce activity duration by 20\%'' (which says \textit{where} to improve but not \textit{how}), we target resource-related improvements: shifting responsibilities between roles, adjusting resource allocation, and modifying working calendars. These intervention types are directly implementable through organizational decisions.

The dual-simulator validation addresses a fundamental challenge in simulation-based assessment: different simulation paradigms capture different aspects of process behavior. DES excels at modeling queue-based capacity effects and resource contention, while ABS captures agent behavioral patterns and handover sequences. By averaging results from both simulators, we obtain validation that is more robust than either alone.

\subsection{Contributions}

This paper makes the following contributions:

\begin{itemize}
    \item An \textbf{agentic AI architecture} for automated process improvement that combines LLM reasoning with simulation-based validation, enabling autonomous exploration of improvement opportunities.

    \item A \textbf{dual-simulator validation methodology} that averages DES and ABS results to provide robust assessment of improvement hypotheses, mitigating the limitations of individual simulation paradigms.

    \item A \textbf{safe intervention palette} focusing on three actionable improvement types (responsibility shifting, resource allocation, calendar adjustment) with role compatibility validation to ensure practical feasibility.

    \item \textbf{Comprehensive evaluation} through quantitative benchmarking against grid search baselines and qualitative assessment with domain experts, demonstrating both efficiency gains and practical utility.
\end{itemize}

The remainder of this paper is organized as follows. Section~\ref{sec:background} provides background on process mining, business process simulation, and LLMs. Section~\ref{sec:related} reviews related work. Section~\ref{sec:framework} presents our framework architecture. Section~\ref{sec:evaluation} describes our experimental evaluation. Section~\ref{sec:threats} discusses threats to validity. Section~\ref{sec:conclusion} concludes and outlines future work.

%% ============================================================================
\section{Background}
\label{sec:background}
%% ============================================================================

This section provides the necessary background on business process management, process mining, simulation approaches, and large language models that underpin our framework.

\subsection{Business Process Management and Event Logs}

Business Process Management (BPM) is a discipline focused on improving organizational performance through the management of business processes~\cite{dumas2018fundamentals}. Modern enterprise information systems record detailed traces of process executions in the form of \textit{event logs}. An event log $L$ consists of a set of traces, where each trace $\sigma = \langle e_1, e_2, \ldots, e_n \rangle$ represents the sequence of events for a single process instance (case). Each event $e$ typically includes attributes such as the activity name, timestamp, and the resource who performed the activity.

Event logs enable data-driven analysis of process performance through metrics such as:
\begin{itemize}
    \item \textbf{Cycle time}: The total duration from case start to completion
    \item \textbf{Waiting time}: Time spent by cases waiting for resources or processing
    \item \textbf{Throughput}: The number of cases completed per unit time
    \item \textbf{Resource utilization}: The fraction of available time that resources spend on productive work
\end{itemize}

\subsection{Process Mining}

Process mining bridges the gap between data science and process management by extracting knowledge from event logs~\cite{van2016process}. The three main types of process mining are:

\begin{itemize}
    \item \textbf{Process Discovery}: Automatically constructing a process model from an event log
    \item \textbf{Conformance Checking}: Comparing observed behavior against a normative model
    \item \textbf{Enhancement}: Extending or improving models based on log analysis
\end{itemize}

Our framework leverages process discovery to extract baseline models and performance enhancement analysis to identify improvement opportunities.

\subsection{Business Process Simulation}

Business Process Simulation (BPS) enables ``what-if'' analysis by executing process models under different scenarios. Two predominant paradigms exist:

\subsubsection{Discrete Event Simulation (DES)}

DES models a process as a sequence of discrete events occurring at specific points in time. Resources are represented as pools with defined capacities, and cases compete for resources through queuing mechanisms. DES is particularly effective at capturing:
\begin{itemize}
    \item Queue-based waiting times due to resource contention
    \item Capacity effects from resource pool sizes
    \item Calendar-based availability constraints
\end{itemize}

However, DES typically does not model individual agent behaviors or complex routing decisions based on agent characteristics.

\subsubsection{Agent-Based Simulation (ABS)}

ABS models processes through autonomous agents that interact according to defined behaviors. Each resource is represented as an individual agent with its own characteristics, decision rules, and performance patterns. ABS excels at capturing:
\begin{itemize}
    \item Individual agent behavioral patterns
    \item Handover sequences and routing preferences
    \item Learning effects and specialization
\end{itemize}

However, ABS typically lacks explicit queue modeling, meaning waiting times arise only from agent availability patterns rather than explicit queuing mechanisms.

\subsubsection{Complementary Strengths}

The complementary nature of DES and ABS motivates our dual-simulator approach. Table~\ref{tab:simulators} summarizes their strengths and limitations.

\begin{table}[htbp]
\centering
\caption{Comparison of DES and ABS simulation paradigms}
\label{tab:simulators}
\begin{tabular}{@{}lll@{}}
\toprule
\textbf{Aspect} & \textbf{DES} & \textbf{ABS} \\
\midrule
Resource modeling & Pools with capacity & Individual agents \\
Queue modeling & Explicit & Implicit/None \\
Waiting time & From queuing & From availability \\
Behavioral patterns & Limited & Rich \\
Handover sequences & Not modeled & Captured \\
Best for validating & Resource allocation, calendars & Responsibility shifts \\
\bottomrule
\end{tabular}
\end{table}

\subsection{Large Language Models and Agentic AI}

Large Language Models are neural networks trained on vast text corpora that demonstrate remarkable capabilities in natural language understanding, reasoning, and generation~\cite{brown2020language}. Recent work has shown that LLMs can be effective at complex planning and decision-making tasks when properly prompted.

\textit{Agentic AI} refers to AI systems that operate with a degree of autonomy, perceiving their environment, reasoning about goals, planning actions, and executing them with minimal human intervention~\cite{wang2024survey}. In industrial contexts, agentic AI enables intelligent automation of complex cognitive tasks that previously required human expertise.

Our framework embodies the agentic paradigm: the LLM autonomously generates improvement hypotheses, the system validates them through simulation, and the iteration continues until a satisfactory improvement is found---all with minimal human intervention during execution.

%% ============================================================================
\section{Related Work}
\label{sec:related}
%% ============================================================================

Our work intersects several research areas: process improvement and prescriptive process mining, simulation-based optimization, LLMs for operations research, and the emerging field of agentic AI for industrial systems.

\subsection{Process Improvement and Prescriptive Analytics}

Traditional process improvement approaches rely on methodologies such as Six Sigma, Lean, and Business Process Reengineering~\cite{dumas2018fundamentals}. These require significant manual effort from domain experts to identify improvement opportunities.

Process mining has evolved from descriptive analytics toward prescriptive capabilities. Bottleneck analysis~\cite{denisov2018predictive} identifies performance constraints, while simulation-based approaches enable ``what-if'' analysis~\cite{camargo2020automated}. SIMOD~\cite{camargo2020automated} automatically discovers simulation models from event logs, enabling data-driven simulation without manual model construction.

However, existing prescriptive approaches typically require either manual hypothesis generation by experts or predefined search spaces for automated optimization. Our work addresses this gap through LLM-based hypothesis generation.

\subsection{Simulation-Based Process Optimization}

Simulation has long been used for process optimization~\cite{law2015simulation}. Traditional approaches formulate optimization problems with explicit objective functions and constraints, then apply techniques such as genetic algorithms or simulated annealing to search for optimal configurations.

Recent work has combined process mining with simulation for improvement analysis~\cite{martin2016recommendations}. However, these approaches typically focus on specific improvement types (e.g., resource allocation) with predefined search spaces. Our framework takes a more general approach, allowing the LLM to propose diverse improvement types while simulation provides validation.

\subsection{LLMs for Operations Research}

The application of LLMs to operations research and optimization problems is an emerging area~\cite{ahmaditeshnizi2024optimus}. Studies have explored using LLMs as optimizers~\cite{yang2024large}, for mathematical modeling~\cite{li2024large}, and for combinatorial optimization~\cite{romera2024mathematical}.

In the process mining domain, LLMs have been applied to process model generation~\cite{berti2023abstractions} and conversational process analysis~\cite{grohs2024large}. However, the combination of LLMs with simulation for iterative process improvement remains largely unexplored.

\subsection{Agentic AI and Digital Twins}

The concept of agentic AI---autonomous systems that perceive, reason, plan, and act---has gained significant attention with the advent of capable LLMs~\cite{wang2024survey}. Multi-agent systems and autonomous decision-making are increasingly applied to industrial contexts.

Digital twins---virtual replicas of physical systems that enable simulation and analysis---are closely related to our work~\cite{tao2018digital}. Our framework can be viewed as creating a ``process digital twin'' from event logs, then using agentic AI to explore improvements on this virtual representation before real-world implementation.

\subsection{Positioning of Our Work}

Our work uniquely combines: (1) LLM-based hypothesis generation for creative exploration, (2) dual-simulator validation for robust assessment, and (3) focus on actionable interventions for practical utility. Unlike pure optimization approaches, we leverage LLM reasoning to propose diverse improvements; unlike pure LLM approaches, we validate all recommendations through rigorous simulation.

%% ============================================================================
\section{Framework Architecture}
\label{sec:framework}
%% ============================================================================

This section presents the architecture of our agentic process improvement framework.\footnote{Source code is available at \url{https://github.com/ANONYMOUS/process-improvement-agent}.} Figure~\ref{fig:architecture} provides an overview of the system.

\begin{figure}[!t]
    \centering
    \includegraphics[width=\textwidth]{architecture.pdf}
    \caption{Framework architecture showing the iterative improvement pipeline}
    \label{fig:architecture}
\end{figure}

\subsection{System Overview}

The framework operates as an iterative pipeline with four main phases:

\begin{enumerate}
    \item \textbf{Baseline Discovery}: Extract a simulation model and baseline performance metrics from the input event log
    \item \textbf{Hypothesis Generation}: Use an LLM to generate improvement hypotheses based on baseline metrics and prior attempts
    \item \textbf{Dual-Simulator Validation}: Execute both DES and ABS with the proposed changes, average results, and assess improvement
    \item \textbf{Report Generation}: Produce actionable recommendations with quantified impacts
\end{enumerate}

Algorithm~\ref{alg:pipeline} presents the main iteration loop.

\begin{algorithm}[t]
\caption{Improvement Validation Pipeline}
\label{alg:pipeline}
\begin{algorithmic}[1]
\REQUIRE Event log $L$, max iterations $N$, improvement threshold $\theta$
\ENSURE First validated improvement $(h^*, score^*)$ or $\emptyset$ if none found
\STATE $M \leftarrow \textsc{DiscoverModel}(L)$ \COMMENT{Baseline model discovery}
\STATE $B_{DES}, B_{ABS} \leftarrow \textsc{SimulateBaseline}(M)$ \COMMENT{Baseline metrics}
\STATE $H_{prev} \leftarrow \emptyset$ \COMMENT{History of attempted hypotheses}
\FOR{$i = 1$ to $N$}
    \STATE $h \leftarrow \textsc{LLMGenerate}(B_{DES}, B_{ABS}, H_{prev})$ \COMMENT{Generate hypothesis}
    \STATE $M' \leftarrow \textsc{ApplyModification}(M, h)$ \COMMENT{Apply to model}
    \STATE $S_{DES} \leftarrow \textsc{SimulateDES}(M')$ \COMMENT{DES simulation}
    \STATE $S_{ABS} \leftarrow \textsc{SimulateABS}(M')$ \COMMENT{ABS simulation}
    \STATE $score \leftarrow \textsc{Validate}(B_{DES}, B_{ABS}, S_{DES}, S_{ABS})$ \COMMENT{Compute score}
    \IF{$score > \theta$ \AND $\neg\textsc{UtilizationVeto}(S_{DES}, S_{ABS})$}
        \RETURN $(h, score)$ \COMMENT{Return first validated improvement}
    \ENDIF
    \STATE $H_{prev} \leftarrow H_{prev} \cup \{(h, score)\}$ \COMMENT{Update history}
\ENDFOR
\RETURN $\emptyset$ \COMMENT{No valid improvement found}
\end{algorithmic}
\end{algorithm}

\subsection{Baseline Model Discovery}

The first phase extracts a simulation model from the event log. We leverage SIMOD~\cite{camargo2020automated}, a state-of-the-art tool that automatically discovers:

\begin{itemize}
    \item A BPMN process model capturing the control flow
    \item Resource profiles defining roles and their capacities
    \item Activity duration distributions fitted to observed data
    \item Resource calendars specifying working hours
    \item Task-resource assignments indicating which resources perform which activities
\end{itemize}

Additionally, we extract baseline Key Performance Indicators (KPIs) directly from the event log using standard process mining techniques:
\begin{itemize}
    \item Average cycle time per case
    \item Throughput (cases per day)
    \item Waiting time between activities
    \item Resource utilization per role
\end{itemize}

These baseline metrics serve dual purposes: guiding the LLM toward high-impact improvements and providing reference points for validation.

\subsubsection{Simulator-Specific Baselines}

Before the improvement iteration loop begins, both simulators execute the discovered model without modifications to establish their own baseline metrics. This dual-baseline approach is essential because DES and ABS may produce different absolute values for the same process due to their different modeling paradigms (e.g., ABS does not model queues, so its waiting time is always zero). By comparing each hypothesis against the corresponding simulator's baseline, we ensure fair evaluation: improvements are measured as relative changes within each simulation paradigm rather than against potentially incompatible absolute values.

\subsection{Hypothesis Generation}

The hypothesis generation component uses an LLM to propose process improvements.\footnote{All prompts used for hypothesis generation and report generation are provided in the repository.} The LLM receives a structured prompt containing:

\begin{itemize}
    \item \textbf{Baseline metrics}: Current performance KPIs and resource utilization
    \item \textbf{Process structure}: Available roles, their tasks, and relationships
    \item \textbf{Improvement palette}: Allowed intervention types and their parameters
    \item \textbf{History}: Previously attempted hypotheses and their outcomes
    \item \textbf{Strategy hints}: Guidance toward high-probability improvements (e.g., targeting high-utilization bottlenecks)
\end{itemize}

The LLM generates a structured hypothesis specifying:
\begin{itemize}
    \item The improvement type
    \item Target role(s) or resource(s)
    \item Specific parameter changes
    \item Expected impact rationale
\end{itemize}

Figure~\ref{fig:prompt} presents the system prompt used for hypothesis generation.

\begin{figure}[!t]
\begin{tcolorbox}[
    title=System Prompt: Hypothesis Generation,
    colframe=black!80,
    fonttitle=\bfseries,
    boxsep=2pt, left=5pt, right=5pt, top=1mm, bottom=1mm,
    toptitle=0.5mm, bottomtitle=0.5mm
]
    \footnotesize
    \begin{tcolorbox}[colback=rolecolor, colframe=white, boxrule=0pt,
        size=fbox, arc=0mm,
        before skip=0pt, after skip=1pt]
        \textbf{Role:} Process improvement specialist analyzing business process metrics.
    \end{tcolorbox}

    \begin{tcolorbox}[colback=goalcolor, colframe=white, boxrule=0pt,
        size=fbox, arc=0mm,
        before skip=0pt, after skip=1pt]
        \textbf{Task:} Generate actionable, predictive hypotheses to improve process performance.
    \end{tcolorbox}

    \begin{tcolorbox}[colback=rulecolor, colframe=white, boxrule=0pt,
        size=fbox, arc=0mm,
        before skip=0pt, after skip=1pt]
        \textbf{Validation Oracle (by change type):}
        \begin{itemize}[nosep, leftmargin=*, labelsep=1.5pt]
            \item \texttt{shift\_responsibility}: Validated by ABS (agent handover). Primary KPI: cycle\_time.
            \item \texttt{resource\_allocation}: Validated by DES (queue model). Primary KPIs: waiting\_time, throughput.
            \item \texttt{calendar\_adjustment}: Validated by DES. Primary KPI: waiting\_time.
        \end{itemize}

        \textbf{Safe Intervention Palette (strict priority):}
        \begin{enumerate}[nosep, leftmargin=*, labelsep=1.5pt]
            \item \texttt{shift\_responsibility} -- Mandatory first choice when overloaded + underutilized roles coexist. Lowest cost: redistributes existing workforce.
            \item \texttt{resource\_allocation} -- Second choice: add capacity to bottleneck roles.
            \item \texttt{calendar\_adjustment} -- Last resort: extend working hours.
        \end{enumerate}

        \textbf{Role Compatibility (for shift\_responsibility):}
        \begin{itemize}[nosep, leftmargin=*, labelsep=1.5pt]
            \item Only shift between roles with overlapping/compatible task assignments.
            \item External actors (e.g., Applicant) cannot perform internal staff tasks.
            \item Verify compatibility using provided role-task mapping.
        \end{itemize}
    \end{tcolorbox}

    \begin{tcolorbox}[colback=outputcolor, colframe=white, boxrule=0pt,
        size=fbox, arc=0mm,
        before skip=0pt, after skip=0pt]
        \textbf{Output:} JSON object with: \texttt{description}, \texttt{improvement\_type}, \texttt{rationale}, \texttt{changes} (parameter\_type, specific\_changes with target/change/value), \texttt{expected\_impact} (cycle\_time, throughput, waiting\_time percentages).
    \end{tcolorbox}
\end{tcolorbox}
\caption{System prompt for LLM-based hypothesis generation (abbreviated)}
\label{fig:prompt}
\end{figure}

\subsubsection{Safe Intervention Palette}

We restrict improvements to three actionable intervention types, prioritized as follows:

\paragraph{Shift Responsibility (Priority 1)}
Redistributing workload between roles that share compatible tasks. This intervention has zero hiring cost and addresses workload imbalances. A shift is validated for role compatibility:
\begin{itemize}
    \item Source and destination must not be external actors (e.g., customers)
    \item Roles must share overlapping tasks in their assignments
    \item Roles must be compatible types (e.g., same role category)
\end{itemize}

\paragraph{Resource Allocation (Priority 2)}
Adding resources to bottleneck roles. This addresses capacity constraints but involves hiring costs. The LLM proposes additions based on utilization levels:
\begin{itemize}
    \item High utilization ($\geq 90\%$): suggest 3--5 additional resources
    \item Moderate utilization ($\geq 80\%$): suggest 2--4 additional resources
\end{itemize}

\paragraph{Calendar Adjustment (Priority 3)}
Extending working hours for constrained resources. This is considered a last resort due to policy and feasibility constraints.

The prioritization reflects both practical feasibility (shifts are easiest to implement) and cost considerations (avoiding unnecessary hiring).

\subsection{Dual-Simulator Validation}

Before validation begins, each simulator computes its baseline metrics by running the unmodified process model over a shared time horizon. This ensures that improvements are measured against comparable reference points, enabling meaningful before-and-after comparisons regardless of the simulator's internal mechanics.

Each hypothesis is then validated through parallel execution of both simulation paradigms.

\subsubsection{DES Validation (Prosimos)}

The DES simulator receives the modified process model and executes discrete event simulation. Modifications are applied to:
\begin{itemize}
    \item Resource profiles (capacity changes for resource allocation)
    \item Resource calendars (hour extensions for calendar adjustment)
    \item Task-resource distribution (reassignments for responsibility shifts)
\end{itemize}

DES is particularly reliable for validating resource allocation and calendar changes due to its explicit queue modeling.

\subsubsection{ABS Validation (AgentSimulator)}

The ABS simulator operates on a modified version of the event log, where:
\begin{itemize}
    \item Resource allocation changes are modeled by distributing cases across synthetic additional agents
    \item Responsibility shifts are modeled by reassigning activities to destination agents
    \item Calendar changes are modeled by compressing activity durations proportionally
\end{itemize}

ABS is particularly reliable for validating responsibility shifts due to its modeling of agent behavioral patterns.

\subsubsection{Validation Formula}

We compute a unified improvement score by averaging simulator results:

\begin{equation}
\label{eq:score}
score = \frac{score_{DES} + score_{ABS}}{2}
\end{equation}

where each simulator's score captures improvement in throughput and cycle time:

\begin{equation}
\label{eq:simulator_score}
score_{sim} = \Delta_{throughput} - \Delta_{cycle\_time}
\end{equation}

Here, $\Delta$ represents the relative change from baseline: $\Delta_m = (m_{new} - m_{baseline}) / m_{baseline}$.

The formulation ensures that both throughput increases (positive $\Delta_{throughput}$) and cycle time decreases (negative $\Delta_{cycle\_time}$) contribute positively to the score.

\subsubsection{Acceptance Criteria}

An improvement hypothesis is accepted if:
\begin{enumerate}
    \item The averaged score exceeds the threshold: $score > 0.01$ (1\% improvement)
    \item The utilization veto is not triggered: average resource utilization $< 90\%$
\end{enumerate}

The utilization veto prevents accepting improvements that overload resources, ensuring practical feasibility.

\subsubsection{Why Dual Validation?}

The rationale for averaging DES and ABS results is threefold:

\begin{enumerate}
    \item \textbf{Complementary coverage}: DES captures queue effects; ABS captures behavioral patterns
    \item \textbf{Bias mitigation}: Averaging reduces simulator-specific biases
    \item \textbf{Robustness}: Improvements must show positive effects in both paradigms
\end{enumerate}

\subsubsection{Why Relative Baselines?}

Simulators may not perfectly match absolute real-world metrics due to modeling assumptions. By comparing improvements against simulator-specific baselines (rather than absolute log metrics), we ensure fair evaluation under consistent assumptions within each simulation paradigm.

\subsection{Report Generation}

The final phase produces an actionable improvement report containing:

\begin{itemize}
    \item \textbf{Executive summary}: Key findings and quantified expected savings
    \item \textbf{Baseline analysis}: Current performance and identified bottlenecks
    \item \textbf{Validated improvement}: Detailed description of the accepted improvement with before/after metrics
    \item \textbf{Implementation roadmap}: Prioritized sequence of changes with dependencies
    \item \textbf{Expected ROI}: Business impact estimates based on simulation results
    \item \textbf{Risks and monitoring}: Metrics to track and warning signs to watch
\end{itemize}

The report is generated using the LLM with strict grounding: all claims must be supported by simulation results, and the narrative explicitly distinguishes between DES and ABS evidence. Figure~\ref{fig:report-prompt} presents the system prompt used for report generation, which enforces factual consistency and prevents hallucination.

\begin{figure}[!t]
\begin{tcolorbox}[
    title=System Prompt: Report Generation,
    colframe=black!80,
    fonttitle=\bfseries,
    boxsep=2pt, left=5pt, right=5pt, top=1mm, bottom=1mm,
    toptitle=0.5mm, bottomtitle=0.5mm
]
    \footnotesize
    \begin{tcolorbox}[colback=rolecolor, colframe=white, boxrule=0pt,
        size=fbox, arc=0mm,
        before skip=0pt, after skip=1pt]
        \textbf{Role:} Business Process Management consultant preparing an executive report for a VP of Operations.
    \end{tcolorbox}

    \begin{tcolorbox}[colback=goalcolor, colframe=white, boxrule=0pt,
        size=fbox, arc=0mm,
        before skip=0pt, after skip=1pt]
        \textbf{Task:} Create a clear, actionable report highlighting validated improvements and their business impact.
    \end{tcolorbox}

    \begin{tcolorbox}[colback=rulecolor, colframe=white, boxrule=0pt,
        size=fbox, arc=0mm,
        before skip=0pt, after skip=1pt]
        \textbf{Grounding Rules:}
        \begin{itemize}[nosep, leftmargin=*, labelsep=1.5pt]
            \item Use ONLY facts present in the provided input payload.
            \item Do NOT invent iterations, targets, people, or metric values.
            \item Keep DES and ABS evidence separate; never merge into synthetic metrics.
            \item If DES and ABS differ, acknowledge divergence and recommend caution.
            \item Do not claim causality beyond what metrics support.
        \end{itemize}
    \end{tcolorbox}

    \begin{tcolorbox}[colback=outputcolor, colframe=white, boxrule=0pt,
        size=fbox, arc=0mm,
        before skip=0pt, after skip=0pt]
        \textbf{Output Structure:} Executive summary (key findings, confidence level), Baseline analysis (bottlenecks, root causes), Validated improvements (before/after metrics, prerequisites, risks), Implementation roadmap (prioritized changes, dependencies), Expected ROI (quantified savings, payback period), Risks \& monitoring (metrics to track, rollback criteria).
    \end{tcolorbox}
\end{tcolorbox}
\caption{System prompt for LLM-based report generation (abbreviated)}
\label{fig:report-prompt}
\end{figure}

%% ============================================================================
\section{Experimental Evaluation}
\label{sec:evaluation}

We evaluate our framework through two complementary studies: a \textit{quantitative} evaluation comparing the recommendation capabilities of different LLMs employed in the framework and how the LLM-guided search for improvements performs against a systematic baseline (i.e., a traditional grid search) with respect to each kind of supported improvement, and a \textit{qualitative} assessment of the actionability and quality of improvement recommendations with process analysts at SAP.

\subsection{Event Logs}
\label{sec:logs}
For the quantitative evaluation, we considered well-established event logs for business process simulation tasks from the process mining community\footnote{All the datasets are available at: \url{}}, listed alongside their properties in Table~\ref{tbl:log_properties}. In particular, we used \textit{BPI Challenge 2012} (work items variant), \textit{BPI Challenge 2017} (work items variant), \textit{Environmental Permit}, \textit{Loan Application}, \textit{Procure-to-Pay}, \textit{Helpdesk}, and a proprietary real-world \textit{Order-to-Cash} (O2C) event log recorded by SAP systems. This selection provides a wide variety of domains and log complexities to assess the generalization capabilities of our LLM agent. Moreover, the inclusion of the proprietary O2C log, which is not contained in the training data of the evaluated LLMs, helps avoid \textit{data contamination} and suggests that results on it reflect real performance rather than memorization of training data~\citep{DBLP:conf/uss/CarliniTWJHLRBS21}.

Specifically, the \textit{BPI Challenge 2012} (BPIC12W) log captures the BP of handling personal loan applications at a Dutch financial institution~\citep{https://doi.org/10.4121/uuid:3926db30-f712-4394-aebc-75976070e91f}. We use the work items (W) variant of the dataset, which focuses on the actual work performed by employees. Despite having few activity types, the BP involves a large resource pool, creating scenarios with parallelism and complex decision logic around checks, approvals, and offers.

Similar to BPIC12W, we consider the work items version of the \textit{BPI Challenge 2017} (BPIC17W) dataset, recording the execution of another loan application BP, capturing the end-to-end workflow from application submission through offer handling~\citep{https://doi.org/10.4121/14625948.v1}. It has the highest number of traces and events in our benchmark, featuring a very large resource pool.

The \textit{Environmental Permit} (Env Permit) dataset originates from the CoSeLoG project and records a building permit application process at a Dutch municipality~\citep{https://doi.org/10.4121/uuid:c45dcbe9-557b-43ca-b6d0-10561e13dcb5}. It has a high activity-to-resource ratio, indicating specialized roles, and exhibits complex handover patterns between organizational units.

The \textit{Helpdesk} log records the ticket handling process of an Italian software company~\citep{https://doi.org/10.4121/uuid:0c60edf1-6f83-4e75-9367-4c63b3e9d5bb}. It captures the lifecycle of support tickets from creation through resolution, including assignment, investigation, escalation, and closure, and exhibits typical service desk patterns with frequent reassignments and escalations.

Regarding synthetically generated logs, \textit{Loan Application} is a dataset recording a loan application process that covers application submission, credit checks, risk assessment, and approval decisions~\citep{https://doi.org/10.4121/uuid:cdf3ba31-291d-468d-9712-3a58ac6da3fc}. Instead, the \textit{Procure-to-Pay} (P2P) dataset is a synthetic event log simulating an end-to-end procurement BP from purchase requisition through invoice payment, spanning from requisition creation to purchase order handling, goods receipt, and invoice processing. The underlying BP features multiple decision points and parallel activities.

Lastly, the \textit{O2C} event log captures a real-world O2C process from a specialty chemicals manufacturer, recorded by SAP systems. It covers the workflow from quotation creation through order processing, delivery, and invoicing, with dedicated roles for sales, warehouse operations, billing, and credit management. Specifically, it has by far the highest number of activities in our benchmark, representing industrial-scale BP complexity. We cannot disclose this log due to a non-disclosure agreement, and for the qualitative evaluation with SAP process analysts, we used a properly anonymized version of it.

\begin{table}[htbp]
\caption{Main statistics of the event logs selected for the evaluation}
\label{tbl:log_properties}
\begin{tabular*}{\tblwidth}{@{} L C R R R R @{}}
\toprule
\textbf{Log} & \textbf{Type} & \textbf{\# Traces} & \textbf{\# Events} & \textbf{\# Activities} & \textbf{\# Resources} \\ 
\midrule
BPIC12W & real & 8616 & 59302 & 6 & 52 \\
BPIC17W & real & 30276 & 240854 & 8 & 136 \\
Env Permit & real & 1434 & 8577 & 27 & 48 \\
Helpdesk & real & 4580 & 21348 & 14 & 22 \\
Loan Application & syn & 1000 & 7492 & 12 & 19 \\
P2P & syn & 608 & 9119 & 21 & 27 \\
O2C SAP & real & 16099 & 54363 & 75 & 14 \\
\bottomrule
\end{tabular*}
\end{table}

\subsection{Experimental Settings}
\label{sec:settings}

\subsubsection{Quantitative Evaluation}
\label{sec:quantitative_setup}
The quantitative evaluation we carried out was two-fold. On the one hand, it aimed to verify how the specific language model impacts the framework performance in proposing an actionable improvement for the BP underlying the considered event log in input (i.e., \textit{``LLM Benchmark''}). On the other hand, fixing the LLM, we wanted to evaluate the performance gain with respect to a basic traditional grid search in the three kinds of supported improvement interventions (i.e., \textit{``Intervention Type Benchmark''}).

For both the assessments, to mitigate the inherent statistical variability of LLMs in generating their outputs, even in response to the same prompt, we repeated the evaluation for 5 runs and averaged the results.

The experiments were conducted on a laptop equipped with a Windows 11 operating system, an 11th-gen Intel Core i5 processor, and 16 GB of RAM. The LLMs were invoked via API calls.

\subsubsection{LLM Benchmark}

We benchmarked multiple LLMs as agents implemented within our framework. In particular, we used the SAP Cloud SDK for AI\footnote{\url{https://sap.github.io/ai-sdk/}} for the implementation and developed the framework in Python 3.11, leveraging a large suite of available LLMs. We selected models to have a wide variety of proprietary, large LLMs as well as open-weights ones, with various sizes in terms of parameters and reasoning capabilities.

\begin{table}[htbp]
\centering
\caption{LLM models evaluated. Size estimates for proprietary models are based on publicly available reports as actual sizes are undisclosed.}
\label{tab:models}
\begin{tabular}{@{}llrlcl@{}}
\toprule
\textbf{Provider} & \textbf{Model} & \textbf{Size (B params)} & \textbf{Reasoning} & \textbf{License} \\
\midrule
OpenAI & GPT 4.1 & $\sim$200--1800 & No & Proprietary \\
OpenAI & GPT 5 Nano & $\sim$8 & No & Proprietary \\
OpenAI & GPT 5 Mini & $\sim$8--70 & No & Proprietary \\
OpenAI & GPT 5 & $\sim$1700--5000 & Yes & Proprietary \\
OpenAI & o3 & $\sim$200--300 & Yes & Proprietary \\
\midrule
Anthropic & Claude 4.5 Haiku & $\sim$20 & No & Proprietary \\
Anthropic & Claude 4.5 Sonnet & $\sim$70--175 & Mixed & Proprietary \\
Anthropic & Claude 4.5 Opus & $\sim$300+ & Yes & Proprietary \\
\midrule
Mistral & Small 3.2 Instruct & 23 & No & Open weights \\
Mistral & Large 3 Instruct & 675 & No & Open weights \\
\bottomrule
\end{tabular}
\end{table}

The metrics we considered for this quantitative assessment are the following:

\begin{itemize}
    \item the \textit{time} (in seconds) needed for the simulations and improvement generation until the first validated improvement;
    \item the number of framework \textit{iterations} (i.e., simulation rounds) until the first validated hypothesis;
    \item the \textit{averaged validation score} according to the validation logic introduced in Equation \ref{eq:score};
    \item the \textit{success rate}, i.e., the percentage of runs finding validated improvement
\end{itemize}

For this evaluation, we considered the BPIC12W, BPIC17W, Loan Application, and the SAP O2C log, to have a comprehensive overview of the framework performance across event logs with different complexity levels.

\subsubsection{Intervention Type Benchmark}
To evaluate the ability to foresee the right intervention for a specific type of improvement, we designed a fault-injection methodology to create controlled bottlenecks and inefficiencies in the BP execution, and assessed the capability of the agentic framework to identify the right intervention to repair it. We compared it in terms of number of iterations (i.e., simulation invocations) with a systematic grid search over the various parametric configurations for the specific improvement task at hand. This was done to reduce the search space, particularly for the grid search, and have a fair comparison.

In particular, the fault injection mechanism first discovers the baseline simulation model from the event log under analysis using SIMOD. Then, it injects a fault to create a degraded simulation model (e.g., reducing a resource profile from 8 to 1 unit), reflects the modification on a faulty version of the event log, and runs the LLM agent to identify the proper intervention to recover from the fault. Its results are then compared with the grid search trying all possible configurations for the intervention type under analysis in terms of the number of iterations needed to converge to the requested improvement.

We designed 15 experiments across the three improvement types:
Regarding the \textit{resource allocation} experiments, we mainly considered the Loan Application event log; an example of injected fault is reducing the number of clerks from 8 to 1.
Instead, concerning the \textit{calendar adjustment} experiments, we considered the Loan Application, P2P, Helpdesk, and Environmental Permit event logs. For instance, a fault of this kind consists in reducing the working calendar of the resource \textit{``clerk''} from 8h to 2h per day.
Finally, for the experiments testing the shift responsibility interventions, we considered the P2P log as it is the most suited for this, featuring the highest number of resources with compatible roles and responsibilities. A fault instance in this case is to move tasks from the resource \textit{``Magdalena Predutta''} to \textit{``Karel de Groot''} to make the former work under stress and make the latter underutilized.

The metric we employed for this part of the quantitative evaluation is the number of iterations or simulation calls needed by both the agentic framework and the grid search to fix the injected fault and restore the baseline simulation model.

\subsubsection{Qualitative Evaluation}
\label{sec:qualitative_setup}

While quantitative metrics demonstrate the numeric performance of the agentic framework, assessing its practical actionability and utility requires a validation with human users. Thus, to evaluate its potential to foster effective decision support, we designed a qualitative study involving six industrial process analysts. Following contemporary methodological consensus \citep{elmholdt2026many}, we prioritized the depth of empirical exposure to process analysis and domain expertise over a large sample size, given the technical nature of validating the answers coming from the agent.

For this, we deliberately did not consider standard academic benchmarks (such as the BPI challenges) in favor of an anonymized version of the real-world SAP O2C event log to ensure that our LLM agent is tested against complex, noisy, and realistic-scale data.

While anonymization was strictly necessary to comply with corporate data privacy policies, care was taken to preserve the structural, temporal, and resource-related semantics (e.g., resource utilization percentages, queueing times, and bottlenecks) required by both the LLM and the human experts to reason about the underlying BP. This choice prevents the ``toy dataset bias'' and forces the experts to evaluate the framework in a realistic, high-stakes scenario.

In particular, we conducted semi-structured expert interviews (approximately 30 minutes each) with professional process analysts and domain experts from SAP Signavio. The sessions were guided using a loosely structured visual workspace to ease cognitive load and elicit more candid feedback.

The interview protocol was divided into four main phases:

\begin{enumerate}
    \item \textit{Participant profiling:} We collected background information, including current role, years of experience with business processes, and familiarity with both O2C BPs and process mining techniques.
    
    \item \textit{Human baseline (blind) assessment:} Participants were presented exclusively with the baseline metrics of the O2C log (e.g., cycle time, throughput, waiting times, and localized resource utilization such as highly utilized Warehouse Staff versus underutilized Sales Representatives). They were then asked two baseline questions: \textit{``Based on this data, what inefficiencies would you expect to find?''} and \textit{``What would be your first recommendation to improve this process?''}.
    
    \item \textit{Agentic framework answer analysis:} The interviewer then revealed the full process improvement report generated by the framework, including the LLM's baseline diagnosis, the validated intervention, and the implementation considerations in terms of risks and metrics to track.
    
    \item \textit{Expert assessment:} Finally, participants evaluated the AI-generated report against their own initial hypotheses.
\end{enumerate}

During the final assessment phase, the semi-structured interviews focused on gathering qualitative insights across several dimensions relevant for the integration of the agentic framework in a real product.

\subsection{Evaluation Results}

We now present the results of our empirical evaluation. We first report quantitative findings from the LLM model benchmark (comparing ten models across four event logs), the intervention type benchmark (comparing LLM-guided search against grid search), and an analysis of the value of iterative refinement. We then describe the qualitative evaluation conducted with domain experts.

\subsubsection{Quantitative Evaluation}
\label{sec:quantitative_res}


\paragraph{LLM Model Benchmark.}
Table~\ref{tab:llm_benchmark} summarizes the performance of the ten LLM models across the four event logs. We report averaged results over five runs per model (three for BPIC12W due to computational constraints), with the number of runs completed without hitting the maximum iteration limit (15) indicated as the success rate.

\begin{table*}[!t]
\centering
\caption{LLM benchmark results averaged across four event logs. Success rate indicates the percentage of runs finding a validated improvement within 15 iterations. Time and iterations refer to the first validated improvement. Score denotes the averaged validation score (Equation~\ref{eq:score}).}
\label{tab:llm_benchmark}
\begin{tabular}{@{}llrrrr@{}}
\toprule
\textbf{Provider} & \textbf{Model} & \textbf{Success Rate} & \textbf{Avg Iterations} & \textbf{Avg Time (s)} & \textbf{Avg Score} \\
\midrule
OpenAI & GPT 4.1 & 94\% & 3.8 $\pm$ 3.7 & 520.5 $\pm$ 375.8 & 0.10 $\pm$ 0.11 \\
OpenAI & GPT 5 & 83\%$^\dagger$ & 2.4 $\pm$ 1.7 & 499.6 $\pm$ 308.4 & 0.14 $\pm$ 0.15 \\
OpenAI & GPT 5 Mini & 83\% & 4.8 $\pm$ 4.0 & 1440.9 $\pm$ 1481.3 & 0.12 $\pm$ 0.15 \\
OpenAI & GPT 5 Nano & 89\%$^\dagger$ & 3.1 $\pm$ 2.4 & 767.2 $\pm$ 718.9 & 0.12 $\pm$ 0.17 \\
OpenAI & o3 & 83\% & 2.9 $\pm$ 2.5 & 746.3 $\pm$ 632.9 & 0.06 $\pm$ 0.09 \\
\midrule
Anthropic & Claude 4.5 Opus & 89\% & 3.5 $\pm$ 4.0 & 624.1 $\pm$ 668.6 & 0.11 $\pm$ 0.12 \\
Anthropic & Claude 4.5 Sonnet & 89\% & 2.3 $\pm$ 1.9 & 474.8 $\pm$ 428.2 & 0.10 $\pm$ 0.14 \\
Anthropic & Claude 4.5 Haiku & 72\% & 2.4 $\pm$ 1.8 & 445.7 $\pm$ 374.3 & 0.60 $\pm$ 1.74 \\
\midrule
Mistral & Small 3.2 Instruct & 50\% & 2.7 $\pm$ 1.3 & 511.7 $\pm$ 393.6 & 0.08 $\pm$ 0.07 \\
Mistral & Large 3 Instruct & 100\% & 3.1 $\pm$ 2.3 & 474.9 $\pm$ 321.5 & 0.09 $\pm$ 0.10 \\
\bottomrule
\multicolumn{6}{l}{\footnotesize $^\dagger$GPT-5 and GPT-5 Nano experienced simulator timeout errors on BPIC12W, affecting success rate.}
\end{tabular}
\end{table*}

The results reveal several key findings. First, all models achieved high success rates ($\geq$80\%) except for Mistral Small (50\%) and Claude 4.5 Haiku (72\%), which struggled particularly on complex logs like BPIC17W. Second, there is no clear correlation between model size and the number of iterations needed: Claude 4.5 Sonnet, a mid-sized model, converged fastest (2.3 iterations on average), while larger models like GPT-5 and Claude 4.5 Opus required comparable or more iterations. Third, validation scores exhibit high variance across logs, reflecting the heterogeneity of the improvement opportunities present in each dataset (O2C scores are substantially higher than Loan Application scores, for instance).

Interestingly, reasoning-enabled models (GPT-5, o3, Claude 4.5 Opus) did not consistently outperform non-reasoning alternatives: both groups achieved comparable success rates ($\sim$85\%) and iteration counts ($\sim$3). This suggests that extended chain-of-thought reasoning provides limited advantage for this task, where the simulation feedback already guides the search effectively.

We also analyzed \textit{first-try success rate}---the proportion of runs where the LLM identifies a valid improvement on the very first iteration. Claude 4.5 Opus achieved the highest rate (62.5\%), followed by Claude 4.5 Sonnet (50.0\%) and Haiku (46.2\%). In contrast, Mistral Small succeeded on the first try only 22.2\% of the time. This metric highlights substantial differences in initial hypothesis quality across models, even when final success rates are similar.

Examining the dual-simulator validation scores, we observed that ABS (Agent-Based Simulation) consistently validated improvements more strongly than DES (Discrete Event Simulation), particularly on the O2C log where ABS scores were 5--10$\times$ higher. This divergence suggests that resource-related improvements have a more pronounced impact on individual agent behavior and queueing dynamics than on aggregate flow-level metrics.

Finally, comparing open-weights models (Mistral) against proprietary alternatives, the gap is modest: 75\% vs 85\% success rate, with comparable iteration counts. The difference is primarily driven by Mistral Small's lower reliability; Mistral Large alone achieves 100\% success rate, demonstrating that open-weights models can match proprietary performance. Claude 4.5 Sonnet emerged as a balanced choice overall, combining 89\% success rate with fast convergence (2.3 iterations) and consistent scores.

A key question is whether the iterative simulation-feedback loop adds value, or if LLMs could identify valid improvements on the first attempt. Figure~\ref{fig:first_vs_best} compares the score of the first proposal generated by each model against the score of the eventually validated improvement. Across all logs and models, the first proposal is typically either rejected (score = 0, marked as infeasible) or achieves a negative or marginal score. For instance, on the O2C log, GPT-5's first proposals average a score of $-0.009$, yet after iterative refinement the validated improvements achieve an average score of $0.33$---a dramatic improvement that demonstrates the necessity of the simulation-in-the-loop approach. Similarly, on BPIC17W, GPT-4.1's first proposals average $-0.046$ before converging to $0.034$ after feedback.

This pattern holds across models and logs: the LLM's initial hypotheses often target suboptimal resources or propose changes of incorrect magnitude. The simulation feedback enables the agent to correct course, ultimately converging to validated improvements that pass both DES and ABS validation.

\begin{figure}
    \centering
    \includegraphics[width=0.95\columnwidth]{first_vs_best_score.pdf}
    \caption{Comparison of first proposal score vs.\ validated improvement score on the SAP O2C log. The first LLM hypothesis is typically rejected or suboptimal; iterative refinement through simulation feedback enables convergence to validated improvements with substantially higher scores. Results averaged over 5 experiments per type.}
    \label{fig:first_vs_best}
\end{figure}

\paragraph{Intervention Type Benchmark.}
Table~\ref{tab:intervention_benchmark} reports the comparison between LLM-guided search and grid search for the fault-injection experiments across the three improvement types.

\begin{table}
\centering
\caption{Comparison of LLM-guided search vs. grid search on fault-injection experiments. Iterations refer to the number of simulation invocations needed to find the corrective intervention. Speedup is computed as grid-search iterations divided by LLM iterations. Results averaged over 5 experiments per type.}
\label{tab:intervention_benchmark}
\begin{tabular}{@{}lrrr@{}}
\toprule
\textbf{Improvement Type} & \textbf{Grid Iters} & \textbf{LLM Iters} & \textbf{Speedup} \\
\midrule
Resource Allocation & 3.6 & 2.0 $\pm$ 0.6 & 1.8$\times$\\
Calendar Adjustment & 4.8 & 1.5 $\pm$ 0.5 & 3.2$\times$ \\
Shift Responsibility & 27.0 & 2.2 $\pm$ 1.0 & 12.3$\times$ \\
\bottomrule
\end{tabular}
\end{table}

The speedup varies substantially by improvement type, reflecting the interplay between search space complexity and the LLM's ability to leverage domain knowledge.

For \textit{resource allocation}, the search space is relatively constrained: given a bottleneck resource, the grid search only needs to enumerate possible capacity increments (e.g., adding 1, 2, or 3 units). In this setting, the LLM provides modest speedup (1.8$\times$) because exhaustive enumeration is already efficient. The LLM's value here lies primarily in identifying \textit{which} resource to target, rather than \textit{how much} to adjust---a decision that both approaches reach similarly quickly once the bottleneck is apparent from the baseline metrics.

For \textit{calendar adjustments}, the benefit increases to 3.2$\times$. Here, the search space expands along multiple dimensions: which resource's calendar to modify, which days or hours to adjust, and by how much. The LLM can reason about temporal patterns in the workload (e.g., peak hours, weekend effects) and propose targeted adjustments, whereas grid search must systematically explore each combination. The LLM's ability to interpret utilization patterns and correlate them with calendar constraints enables more directed exploration.

The most dramatic gains occur for \textit{shift responsibility} interventions, where the LLM achieves 12.3$\times$ speedup on average. This improvement type exhibits combinatorial complexity: with $n$ resources, there are $O(n^2)$ possible source-target pairs for responsibility transfers, and each transfer can involve varying fractions of workload. Grid search must enumerate these combinations systematically, requiring 27 iterations on average in our experiments. In contrast, the LLM can directly analyze workload imbalances---identifying overloaded resources (high utilization, long queues) and underutilized counterparts (low utilization, compatible skills)---and propose targeted transfers in just 2.2 iterations. This represents a qualitative shift from blind enumeration to informed reasoning.

These results have practical implications for framework deployment. When the improvement type is known \textit{a priori} and the search space is small, simpler optimization methods may suffice. However, in realistic scenarios where the improvement type is unknown or where multiple intervention types must be explored simultaneously, the LLM's ability to reason across intervention categories and prioritize promising directions becomes increasingly valuable. The variance in LLM iterations ($\pm$0.5 to $\pm$1.0) also suggests that the approach is reasonably stable, though practitioners should expect some variability in convergence time.

A limitation of this comparison is that grid search serves as a lower bound on systematic exploration rather than a competitive baseline. More sophisticated search strategies (e.g., Bayesian optimization, genetic algorithms) might narrow the gap, particularly for resource allocation and calendar adjustments. However, such methods would require problem-specific encodings and hyperparameter tuning, whereas the LLM operates directly from natural language descriptions of the process and its constraints.


\subsubsection{Qualitative Evaluation}
\label{sec:qualitative_res}

\subsubsection{Study Design}

To assess the practical utility of generated recommendations, we conducted a qualitative study with domain experts. The study aimed to answer: ``Do domain experts find the recommendations actionable and useful?''

\subsubsection{Participants}

We recruited process analysts and domain experts with experience in business process improvement. Participants were presented with improvement recommendations generated by our framework on a real-world process log.

\subsubsection{Procedure}

\begin{enumerate}
    \item Participants reviewed the baseline process analysis
    \item Participants examined each generated improvement recommendation
    \item Participants rated the actionability and usefulness of recommendations
    \item Semi-structured interviews explored reasoning and concerns
\end{enumerate}

\subsubsection{Findings}

% TODO: Add actual qualitative findings
% Placeholder for qualitative results

Domain experts generally found the recommendations actionable, particularly appreciating that improvements targeted specific resource-related interventions rather than abstract ``reduce duration'' suggestions. Experts noted that the dual-simulator validation provided confidence in the predicted impacts.

Key feedback themes included:
\begin{itemize}
    \item Appreciation for the focus on implementable changes
    \item Interest in cost-benefit analysis integration
    \item Requests for more context about why certain roles were targeted
    \item Suggestions for incorporating organizational constraints
\end{itemize}

%% ============================================================================
\section{Threats to Validity}
\label{sec:threats}
%% ============================================================================

We discuss threats to validity organized by category.

\subsection{Internal Validity}

\begin{table}[htbp]
\centering
\caption{Internal validity threats and mitigations}
\label{tab:internal}
\begin{tabular}{@{}p{4cm}p{8cm}@{}}
\toprule
\textbf{Threat} & \textbf{Mitigation} \\
\midrule
LLM stochasticity & Temperature set to 0.0; 5 repetitions per experiment; report mean $\pm$ std \\
Simulation variance & Average 5--10 DES/ABS runs per evaluation \\
Unfair comparison & Simulation cache ensures same inputs produce same outputs \\
Single-simulator bias & Validate with both DES and ABS for all improvement types \\
\bottomrule
\end{tabular}
\end{table}

\subsection{External Validity}

\begin{table}[htbp]
\centering
\caption{External validity threats and mitigations}
\label{tab:external}
\begin{tabular}{@{}p{4cm}p{8cm}@{}}
\toprule
\textbf{Threat} & \textbf{Mitigation} \\
\midrule
Limited log diversity & Use logs from 3 different process domains \\
Artificial fault injection & Inject realistic faults representing bottleneck scenarios \\
Single LLM configuration & Benchmark multiple models; report comparative results \\
\bottomrule
\end{tabular}
\end{table}

\subsection{Construct Validity}

\begin{table}[htbp]
\centering
\caption{Construct validity threats and mitigations}
\label{tab:construct}
\begin{tabular}{@{}p{4cm}p{8cm}@{}}
\toprule
\textbf{Threat} & \textbf{Mitigation} \\
\midrule
KPI selection & Use established process mining metrics (cycle time, throughput, waiting time) \\
Convergence definition & Multi-KPI criterion with primary KPI and throughput guard \\
Validation logic simplicity & Current formula is intentionally simple; future work will expand to multi-objective optimization \\
\bottomrule
\end{tabular}
\end{table}

\subsection{Qualitative Evaluation Limitations}

The qualitative evaluation faces specific limitations:
\begin{itemize}
    \item \textbf{Subjectivity}: Expert opinions inherently vary
    \item \textbf{Expert pool}: Internal recruitment may introduce organizational bias
    \item \textbf{Single log}: Using one log for qualitative assessment may limit generalizability, though this design choice avoided overwhelming participants and maximized feedback quality within limited interview time
\end{itemize}

%% ============================================================================
\section{Discussion}
\label{sec:discussion}
%% ============================================================================

\subsection{Implications for Practice}

Our framework demonstrates the potential for agentic AI to augment process analysts' capabilities. Rather than replacing human expertise, the system serves as an intelligent assistant that:
\begin{itemize}
    \item Explores improvement opportunities that humans might overlook
    \item Validates hypotheses through rigorous simulation before implementation
    \item Generates actionable recommendations with quantified impacts
\end{itemize}

The focus on actionable interventions (responsibility shifts, resource allocation, calendar changes) ensures that recommendations can be directly implemented through organizational decisions.

\subsection{Implications for Research}

Our work opens several research directions at the intersection of LLMs, simulation, and process mining:
\begin{itemize}
    \item \textbf{Multi-objective optimization}: Extending the validation formula to balance multiple competing objectives
    \item \textbf{Cost-aware improvements}: Incorporating financial costs into the optimization
    \item \textbf{Explanation and trust}: Enhancing the interpretability of LLM reasoning
    \item \textbf{Continuous improvement}: Extending to cumulative recommendation sequences
\end{itemize}

%% ============================================================================
\section{Conclusion and Future Work}
\label{sec:conclusion}
%% ============================================================================

We presented an agentic AI framework for automated process improvement that combines LLM-based hypothesis generation with dual-simulator validation. Our approach addresses a key challenge in prescriptive process mining: generating actionable improvement recommendations that are both creative (leveraging LLM reasoning) and validated (through rigorous simulation).

The framework's key innovations include: (1) using LLMs to autonomously generate diverse improvement hypotheses; (2) validating through both DES and ABS to leverage their complementary strengths; (3) focusing on a safe intervention palette of actionable resource-related changes; and (4) providing robust validation through averaged dual-simulator results.

Our evaluation demonstrates that LLM-guided search achieves significant speedup over systematic grid search while producing improvements that domain experts find actionable and useful.

\subsection{Future Work}

We identify several directions for future research:

\begin{itemize}
    \item \textbf{Enhanced Input Data}: Incorporating SAP configuration options, organizational models, and domain knowledge beyond event logs to better define the space of possible improvements.

    \item \textbf{Data Anonymization}: Addressing the challenge that anonymizing logs (required for privacy) results in loss of semantic information that LLMs need for useful recommendations.

    \item \textbf{Expanded Improvement Scope}: Including activity reordering and parallelization, though these require advances in simulator capabilities.

    \item \textbf{Cost Analysis}: Evaluating whether performance gains justify intervention costs (e.g., hiring additional resources).

    \item \textbf{Case Complexity}: Accounting for case features (e.g., claim complexity) to avoid misleading conclusions about resource performance.

    \item \textbf{Mixed Improvement Types}: Enabling the LLM to freely choose among all improvement types in a single search, with appropriate baseline comparison methodology.

    \item \textbf{Cumulative Recommendations}: Testing system behavior across sequences of improvements rather than single recommendations.

    \item \textbf{Goal-Oriented Optimization}: Extending validation to support explicit user-defined optimization goals and constraints.
\end{itemize}

\section*{Declaration of Competing Interest}

The authors declare that they have no known competing financial interests or personal relationships that could have appeared to influence the work reported in this paper.

\section*{Acknowledgments}

% TODO: Add acknowledgments

\bibliographystyle{elsarticle-num}
\bibliography{references}

\end{document}
